\scanpage{106}
Soient $G\longrightarrow G'\simeq G/\unsure{C}$ un hom. de générisation de graphes MD, où $C\subset A=A(G)$ (ens. des arêtes \unsure{liées|libres} de $G$). Soit $X$ une courbe algébrique sur $\mathbb{C}$, de type $G$. Je vais décrire une méthode pour construire des courbes \inserted{$X'$} de type $G'$. On a
\[
C\subset A=\mathrm{Sing}(X),\quad\text{soit }B=A\setminus C
\]
Soit \deleted{\unsure{C}} $\widetilde X$ le normalisé de $X$ relativement à $C$, $\widetilde C$ \inserted{et $\widetilde B$} les images inverses de $C$ \inserted{et $B$} dans $\widetilde X$. Choisissons pour tout \deleted{\unsure{$\widetilde x$}} $x\in\widetilde C$ un disque fermé $D_x\subset\widetilde X_{\deleted{\unsure{\mathrm{sing}}}}\setminus\widetilde B$, contenant $x$ dans son intérieur, à bord \deleted{\unsure{$\mathbb{C}$}} \deleted{tel que les} une courbe analytique réelle (\unsure{pour être strict} --- $C^1$ aurait suffi), de façon que les $D_x$ soient mutuellement disjoints. Il doit être vrai que l'ensemble des choix possibles est \deleted{\unsure{donc}} un espace contractile --- i.e. que le choix est ``anodin'' : la projection de la multiplicité \inserted{top.} formée des $X$ de type $G$, \uline{munis} d'une famille $D_x$ (--- \unsure{j'appeller.} des $X$ $C$-découpés) vers celle des $X$ de type $G$ sans plus, doit être une équivalence d'homotopie. Soit
\[
X_1^{\unsure{\bullet|\circ}}=\widetilde X-\bigcup_{x\in\widetilde C}D_x^\circ
\]
qui est, \deleted{en dehors de \unsure{$X_1\cap\mathrm{Sing}(\widetilde X)$},} une surface à bord holomorphe, avec
\[
\partial X_1=\coprod_{x\in\widetilde C}\partial D_x
\]
En tant que bord d'un disque holomorphe $D_x$, de centre $x$, $\partial D_x$ hérite une structure de $\mathbb{U}$-torseur,

\mspage{105}
